% ================================================================
% 00_frontmatter/foreword_and_preface.tex — लेखक-वक्तव्य आ प्राक्कथन
% ================================================================

% ===================== लेखक के ओर से (Preface) =====================
\cleardoublepage
\thispagestyle{plain}
\vspace*{\fill}
\begin{center}
  {\Large \bfseries \color{bordercolor} लेखक के ओर से} \\[0.4em]
  {\color{bordercolor}\rule{0.25\textwidth}{0.8pt}}
\end{center}
\phantomsection\addcontentsline{toc}{section}{लेखक के ओर से}

\vspace{0.8em}
संत दरिया बाबा के जीवन आ बानी के एह संग्रह के भोजपुरी अनुवाद प्रस्तुत करत हम बड़ा संतोष महसूस करत बानी। मूल कृति एन.के. उपाध्याय द्वारा लिखल गइल रहे, आ एह में दरिया बाबा के अमर उपदेशन के संकलन बा।

ई पुस्तक तीन भाग में बँटल बा — पहिला भाग में जीवन आ उपदेश के विवरण, दुसरका भाग में चुने हुए पद आ बानी के ग्यारह खण्ड में वर्गीकरण, आ अंत में शब्द-सूची।

एह अनुवाद में प्रामाणिक भोजपुरी के प्रयोग कइल गइल बा ताकि दरिया बाबा के बानी आम जनता तक सरलता से पहुँच सके।

\begin{flushright}
— {\bfseries आदित्य कुमार}
\end{flushright}
\vspace*{\fill}



% ===================== प्राक्कथन (Foreword) =====================
\cleardoublepage
\section*{प्राक्कथन}
\phantomsection\addcontentsline{toc}{section}{प्राक्कथन}

\begin{quote}
हे सूरमा! तू काल के नगरी में जा,\\
सुत्तल जीव सब के जगावा आ ओहन के परम आनंद के सागर में वापस ले आव।\\
जब तू माता के गर्भ में प्रवेश करके मनुष्य रूप में अवतार लेब,\\
तबही एह संसार के लोग तोहरा माध्यम से हमार भीतरी गुप्त रहस्य के जान पाई॥\footnote{ज्ञान दीपक, चौपाई ७५७--७५८, मूल पांडुलिपि, पृ. १६७}
\end{quote}

एह बचन के साथ, दरिया बाबा के अनुसार, सत्य पुरुष आपन पुत्र (सतगुरु) के एह संसार में भेजलें। दया के सागर परमेश्वर अनादि काल से आपन दया के काज पर बार-बार आपन संतन के एह जग में भेजत बाड़ें, ताकि जीव सब के काल के पंजा से बचा के ओहन के असली आनंदमय घर (अमर लोक) में वापस पहुँचावल जा सके। बाकिर ई अंधलोक संतन के महिमा समझे के बजाय ओहन के अनदेखा करे ला, मज़ाक उड़ावे ला, अपमानित करे ला आ सतावे ला। संत दरिया बाबा जइसन उच्च कोटि के महापुरुष के साथ भी ऐसने बर्ताव भइल, जइसन कि एह पुस्तक के पहिला खंड — जीवन — से साफ़ पता चले ला।

आज उनकर शरीर छूटे के दो सौ बरिस से ढेर बीत गइला के बाद भी, ऊ आपन जिला या राज्य में भी बहुत कम जानल जात रहलें, देश आ दुनिया के त कोन कहें। बाकिर उनकर उपदेश, जे कवनो धार्मिक हठधर्मिता आ पूर्वाग्रह, संप्रदायवादी अंधविश्वास आ संकीर्णता से मुक्त बा, सार्वभौमिक बा आ पूरा दुनिया के सच्चा साधक आ भक्तन खातिर बा। एह समय तक, दरिया के जीवन आ उपदेश पर अंग्रेजी में कवनो पुस्तक उपलब्ध नइखे, सिवाय फ्रांसिस बुकानन के शाहाबाद रिपोर्ट\footnote{फ्रांसिस बुकानन सन् १८०९--१८१० में शाहाबाद जिला के भ्रमण कइलें, यानी दरिया के निधन के मात्र तीस बरिस बाद। अपने रिपोर्ट में दरिया के उत्तराधिकारी आ उनकर संप्रदाय के बारे में बतवलें।} में संक्षिप्त उल्लेख आ बी.बी. मजुमदार के छोटका लेख {\engfont ``The Tailor Saint of Bihar''} जे ११ सितंबर १९३५ के अंग्रेजी अखबार {\engfont \textit{Searchlight}} में छपल रहे। हमरे जानकारी अनुसार, दरिया पर अंग्रेजी में एगो अउरी छोटका संदर्भ पी.डी. बर्थवाल के कृति {\engfont \textit{Traditions of Indian Mysticism Based upon Nirguna School of Hindi Poetry}} में मिले ला, जेह में बहुत कम जानकारी बा।

पहिली बेर, दरिया के दर्शन आ साहित्य पर गंभीर शोध के कोसिस डॉ. धर्मेंद्र ब्रह्मचारी शास्त्री कइलें, जे दरिया के रचना सभ के आधार पर आपन पीएचडी शोध-प्रबंध तइयार कइलें। ई शोध-प्रबंध हिंदी में अनुवादित भइल आ बिहार राष्ट्रभाषा परिषद द्वारा १९५४ में दरिया के कृति के पहिला खंड के रूप में \textit{दरिया ग्रंथावली, भाग १} शीर्षक से प्रकाशित भइल। बाद में, डॉ. शास्त्री द्वारा दरिया के छह गो पुस्तक संकलित क के १९६२ में राष्ट्रभाषा परिषद द्वारा डॉ. शास्त्री के काम के दुसरका खंड के रूप में \textit{दरिया ग्रंथावली, भाग २} शीर्षक से प्रकाशित भइल। एह में दरिया के निम्नलिखित छह गो पुस्तक सामिल बा: (१) \textit{दरिया सागर}, (२) \textit{ज्ञान रतन}, (३) \textit{ज्ञान स्वरोदय}, (४) \textit{भक्ति हेतु}, (५) \textit{ब्रह्म विवेक}, आ (६) \textit{ज्ञान मूल}। एह पुस्तक के दुसरका खंड के प्रस्तावना में, तत्कालीन परिषद निदेशक भुवनेश्वरनाथ मिश्र "माधव" के अनुसार, परिषद द्वारा दरिया के बाकी रचना सभ के प्रकाशित करे के योजना रहे। बाकिर ओह योजना के लागू होखे से पहिले, डॉ. माधव आ डॉ. शास्त्री दुनों के मृत्यु भ गइल। डॉ. माधव के उत्तराधिकारी लोग एह बहुमूल्य परियोजना पर ध्यान देवे के जरूरी ना समझल। चूँकि दरिया के सगरी मूल पांडुलिपि (एक के छोड़ के जे फारसी भाषा में बा) कैथी लिपि में बाड़ी, ओह सब के पहिले देवनागरी लिपि (आधुनिक हिंदी के लिपि) में लिखल जरूरी बा, तब प्रकाशित हो सके ला। कुछ अप्रकाशित रचना सभ डॉ. शास्त्री के निगरानी में पहिले देवनागरी में लिखल जा चुकल रहलीं आ लगभग प्रेस में जाए खातिर तइयार रहलीं। बाकिर डॉ. माधव के मृत्यु के बाद पूरा परियोजना ठप हो गइल, आ एह में से कुछ बहुमूल्य पांडुलिपि अब कीड़ा के खाद्य बन रहल बाड़ी।

देवनागरी लिपि में पहिले लिखल जा चुकल पांडुलिपि सभ के जेरॉक्स प्रति प्राप्त करे में आ कैथी पांडुलिपि सभ के देवनागरी में लिखवावे में हमरा काफी कठिनाई के सामना करे पड़ल। बाकिर हम बिहार राष्ट्रभाषा परिषद के आभारी बानी जे ऊ अंततः हमरा एह पांडुलिपि सभ के उपयोग करे के अनुमति देलक। हम चाहित बानी कि बिहार राष्ट्रभाषा परिषद या धरकंधा मठ के अधिकारी लोग दरिया बाबा के एह अत्यंत बहुमूल्य पांडुलिपि सभ के प्रकाशित करे के उचित समझी।

कुल २१ गो पुस्तक के बारे में हमरा जानकारी भइल, जे में से दू गो — १. \textit{कल चरित} आ २. \textit{ब्रह्म प्रकाश} — के प्रति हम प्राप्त ना कर सकलीं। बाकिर बाकी बचल प्रकाशित-अप्रकाशित पुस्तक (एगो संस्कृत में आ एगो फारसी में) हम एह पुस्तक के तइयारी में उपयोग कइलीं:

१. \textit{अग्र ज्ञान}, २. \textit{अमर सार}, ३. \textit{भक्ति हेतु}, ४. \textit{ब्रह्म चैतन्य} (संस्कृत), ५. \textit{ब्रह्म विवेक}, ६. \textit{दरियानामा} (फारसी), ७. \textit{दरिया सागर}, ८. \textit{गणेश गोष्ठी}, ९. \textit{ज्ञान दीपक}, १०. \textit{ज्ञान मूल}, ११. \textit{ज्ञान रतन}, १२. \textit{ज्ञान स्वरोदय}, १३. \textit{मूर्ति उखाड़}, १४. \textit{निर्भय ज्ञान}, १५. \textit{प्रेम मूल}, १६. \textit{शब्द}, १७. \textit{सहस्राणी}, १८. \textit{विवेक सागर}, आ १९. \textit{यज्ञ समाधि}।

पुस्तक \textit{दरिया सागर} आ दरिया के बिबिध रचना से लिहल छोटका संग्रह \textit{दरिया बाबा के चुने हुए शब्द} — जे बेलवेडियर प्रेस, इलाहाबाद से प्रकाशित भइल — हम उपयोग कइलीं। पुस्तक \textit{ज्ञान दीपक}, जे साधु चतुरी दास द्वारा १९३६ में दरिया संप्रदाय के अनुयायी लोग में बाँटे खातिर प्रकाशित कइल गइल रहे, उपलब्ध ना रहे। बाकिर चूँकि हम पहिले से एह कृति के पांडुलिपि रूप में प्राप्त कर चुकल रहीं, एकर अप्राप्ति से कवनो फरक ना पड़ल।

ऊपर बतावल पुस्तक आ पांडुलिपि सभ दरिया के उपदेश के साफ़ आ व्यवस्थित रूप से प्रस्तुत करे खातिर पर्याप्त सामग्री उपलब्ध करवलक। दरिया के उपदेश के संक्षिप्त सारांश जइसन कि एह पुस्तक के दुसरका खंड — \textit{उपदेश} — में दिहल बा, आ एकर अधिक विस्तृत विवरण भाग २: \textit{चुने हुए पद आ बानी} में, साफ़ देखावे ला कि दरिया एगो सिद्ध संत आ निडर मालिक रहलें। ऊ चौथा लोक के दयालु प्रभु के निवास आ आत्मा सभ के मूल घर बतावे लें, जे तीनों लोक से परे बा जहाँ काल ना पहुँच सके। एह सबसे उच्च लोक के दरिया \textit{सत लोक} (सत्य लोक) या \textit{छप लोक} (गुप्त लोक) कहतारे, काल के झूठा आ भ्रामक दुनिया के बिपरीत। कबीर निहुँना, दरिया भी सबसे ऊँच लोक खातिर \textit{अकह} (अकथनीय) शब्द के उपयोग करे लें, जे कबीर आ अउरी संतन के अनुसार \textit{अनामी पुरुष} के निवास ह। दरिया एह प्रभु के \textit{बेबहा} (अत्यंत मूल्यवान) कहतारे। दरिया, एक ओर परम प्रभु के महिमा अत्यंत प्रेम आ नम्रता से गावे लें, त दुसरका ओर धार्मिक हठधर्मिता आ कर्मकांड के कड़वा आ व्यंग्यात्मक शब्दन में आलोचना करे लें। उनकर पूरा रचना सतगुरु, शब्द, सत पुरुष या सत नाम के महिमा आ कर्मकांड के निंदा से भरल बा। आध्यात्मिकता के मूल सत्य हमार मन में बैसावे खातिर, दरिया ओह सब के बिबिध तरीका से बतावे लें, हर बेर नवा स्वाद के साथ। ई कुछ दोहराव जइसन लाग सके ला, बाकिर ई सत्य कई बेर दोहरावे लायक बाड़ें।

चूँकि दरिया के सच्चा रहस्यमयी संदेश पर हिंदी या अंग्रेजी में कवनो अउरी पुस्तक नइखे (डॉ. शास्त्री के पुस्तक, उनकर अत्यंत परिश्रमी आ ईमानदार कोसिस के बावजूद, दरिया के उपदेश में वैष्णववाद, शैववाद, वेदांत आ हठ योग आदि के मिलल-जुलल दर्शन खोजे के प्रवृत्ति रखे ला), हम सोचलीं कि \textit{चुने हुए पद आ बानी} के भाग के और ना छोट करल जाय। दरिया के सच्चा रहस्यमयी संदेश उनकर अपना शब्दन में सबसे साफ़ उभर के आवे ला।

कुछ संक्षिप्तियाँ जे एह पुस्तक में प्रयोग भइल बाड़ी, ओह पर ध्यान देवे लायक बा। दरिया ग्रंथावली के दू गो प्रकाशित खंड सभ के संदर्भ में "द.ग्रं. भाग १ आ भाग २" आ अप्रकाशित हस्तलिखित पांडुलिपि सभ के खाली "पांडु." कहल गइल बा, आगे पुस्तक या पांडुलिपि के पूरा शीर्षक जोड़ के। एगो अउरी बात ध्यान देबे लायक बा कि कुछ चुनल पद सभ के संदर्भ खाली पन्ना नंबर से दिहल गइल बा जबकि दोसरा के पन्ना नंबर आ पद नंबर दुनों से। एह से कारण ई बा कि पहिला मामिला में पांडुलिपि में पद नंबर ना दिहल गइल बाड़ें जबकि दुसरका में पद नंबर पांडुलिपि में मौजूद बाड़ें।

अंत में, हम प्रोफ़ेसर के.एस. नारंग, पूर्व कुलपति, पंजाबी यूनिवर्सिटी, पटियाला, भारत के प्रति आभार प्रकट करित बानी, जे हमरा एह परियोजना पर काम करे के सौभाग्य देलें। श्रीमती डायेन ग्रुएनस्टीन, मेयो हैरिस, श्रीमती हेलन सिल्वरबर्ग, श्री एड केनेडी आ उनकर पत्नी श्रीमती मार्गारीटा केनेडी, सब हवाई, अमेरिका के, आ कमला समतानी, राधा स्वामी कॉलोनी, बियास, भारत के, हम ओह सब के आभारी बानी जे पांडुलिपि के बिबिध चरणन में टाइप करे में सहायता कइलीं। हमरा हार्दिक धन्यवाद श्री डेविड सैनफोर्ड आ उनकर पत्नी श्रीमती नोर्मा सैनफोर्ड के भी बा, जे पांडुलिपि के ध्यान से पढ़ के मूल्यवान सुधार कइलें।

\begin{flushright}
यूनिवर्सिटी ऑफ़ हवाई\\
होनोलुलु, हवाई, अमेरिका\\
८ अक्टूबर १९८६

{\bfseries के.एन. उपाध्याय}\\
{\small प्रोफ़ेसर, दर्शनशास्त्र विभाग}
\end{flushright}
